% !TeX encoding = UTF-8
\documentclass[variant=guia,cover=institutional,mode=screen]{ua-engineering-v6-article}
% !TeX encoding = UTF-8
\UASetUniversity{Universidad Austral}
\UASetFaculty{Facultad de Ingeniería}
\UASetDepartment{Departamento de Computación}
\UASetCampus{Pilar}
\UASetCareer{Ingeniería Informática}
\UASetCommission{Comisión A}
\UASetProfessor{Equipo docente}
\UASetSemester{Segundo cuatrimestre 2026}
\UASetCourse{Sistemas Distribuidos}
\UASetDate{2026}

\UASetClassNumber{09}
\UASetClassTitle{Arquitecturas distribuidas reales}
\title{Clase 09 --- Guía de ejercicios}
\author{\UAProfessor}
\date{\UADate}
\begin{document}
\maketitle
\section{Propósito de la guía}
Los ejercicios se evalúan por la coherencia entre dominio, ownership, comunicación, datos y operación. No por cantidad de servicios.
\section{Ejercicios}
\begin{uaexercise}
\textbf{1. Monolito modular.} Explique por qué una unidad de despliegue no implica ausencia de modularidad. Proponga reglas internas para evitar una bola de barro.
\end{uaexercise}
\begin{uaexercise}
\textbf{2. Prima de microservicios.} Enumere los costos distribuidos que aparecen al extraer un módulo y diseñe un criterio para decidir si el beneficio compensa.
\end{uaexercise}
\begin{uaexercise}
\textbf{3. Señales de separación.} Dado un sistema con catálogo, pagos y recomendaciones, decida qué separar primero y qué mantener junto. Justifique por dominio, carga y equipo.
\end{uaexercise}
\begin{uaexercise}
\textbf{4. Monolito distribuido.} Construya una arquitectura con cinco servicios que conserve despliegue coordinado, base compartida y cadenas RPC. Explique por qué es peor que un monolito modular.
\end{uaexercise}
\begin{uaexercise}
\textbf{5. Bounded contexts.} Defina tres significados distintos de 'curso' o 'pedido' en contextos de negocio y diseñe traducciones entre ellos.
\end{uaexercise}
\begin{uaexercise}
\textbf{6. Ownership de datos.} Explique qué significa database per service y cómo resolvería una consulta global sin permitir acceso directo a tablas ajenas.
\end{uaexercise}
\begin{uaexercise}
\textbf{7. Conway.} Analice una organización donde dos equipos deben aprobar cada cambio del otro. ¿Qué revela sobre la frontera de servicios?
\end{uaexercise}
\begin{uaexercise}
\textbf{8. RPC vs eventos.} Para validación de pago, notificación, actualización de ranking y consulta de perfil, elija RPC o evento y justifique.
\end{uaexercise}
\begin{uaexercise}
\textbf{9. Cadena sincrónica.} Diseñe una cadena API->A->B->C y calcule disponibilidad compuesta si cada servicio tiene 99.9%. Analice además latencia y timeouts.
\end{uaexercise}
\begin{uaexercise}
\textbf{10. Comando, evento y consulta.} Clasifique ReserveSeat, SeatReserved y GetSeatAvailability. Construya un ejemplo de evento mal modelado como comando oculto.
\end{uaexercise}
\begin{uaexercise}
\textbf{11. Dual write.} Un servicio actualiza Enrollment y publica StudentEnrolled. Dibuje las ventanas de falla para ambos órdenes.
\end{uaexercise}
\begin{uaexercise}
\textbf{12. Transactional outbox.} Diseñe schema mínimo, relay, deduplicación y política de retención para una outbox de inscripciones.
\end{uaexercise}
\begin{uaexercise}
\textbf{13. Evolución de eventos.} Proponga una evolución compatible de StudentEnrolled v1 a v2 y una evolución que rompería consumidores.
\end{uaexercise}
\begin{uaexercise}
\textbf{14. Saga de inscripción.} Diseñe pasos, timeouts, reintentos y compensaciones para inscripción, cupo, pago, acceso y comprobante.
\end{uaexercise}
\begin{uaexercise}
\textbf{15. Compensación imperfecta.} Dé tres efectos que no pueden revertirse perfectamente y proponga una compensación de negocio razonable.
\end{uaexercise}
\begin{uaexercise}
\textbf{16. Coreografía.} Diseñe una saga coreografiada con tres participantes y explique cómo evitar ciclos y pérdida de visibilidad.
\end{uaexercise}
\begin{uaexercise}
\textbf{17. Orquestación.} Diseñe el mismo flujo con un orquestador durable. Defina estado, recuperación y riesgo de centralización.
\end{uaexercise}
\begin{uaexercise}
\textbf{18. Strangler Fig.} Proponga una migración incremental desde un monolito de e-commerce. Defina primer servicio, routing, rollback y métricas.
\end{uaexercise}
\begin{uaexercise}
\textbf{19. Capacidad operativa.} Construya una checklist de capacidades mínimas antes de operar 30 microservicios en producción.
\end{uaexercise}
\begin{uaexercise}
\textbf{20. Integración D=(P,F,G,S,T).} Aplique el marco a una plataforma de cursos y defienda qué permanece monolítico, qué se extrae y qué trade-offs se aceptan.
\end{uaexercise}
\section{Criterios de entrega}
Toda respuesta debe indicar beneficio, prima distribuida, falla nueva, mecanismo de mitigación y condición para revisar la decisión.
\end{document}
